% Slides — Capítulo 16: Ciclones Tropicais
\documentclass[aspectratio=169,11pt]{beamer}
\usepackage{../comum/temameteo}
\graphicspath{{../figuras/cap16/}}

\title[Cap. 16 --- Ciclones Tropicais]{Capítulo 16 --- Ciclones
Tropicais}
\subtitle{Meteorologia Prática}
\author{}
\date{}

\begin{document}

\begin{frame}
  \titlepage
  \vfill
  \atribuicao
\end{frame}

\begin{frame}{Roteiro}
  \tableofcontents
\end{frame}

% ---------------------------------------------------------------
\section{Anatomia e gênese}

\begin{frame}{Outro animal: o núcleo quente}
  \begin{columns}
    \column{0.3\textwidth}
    \begin{itemize}
      \item Olho (subsidência, calmaria), \alert{parede} (os ventos
            máximos), bandas espirais, saída anticiclônica no topo.
      \item Axissimétrico, sem frentes: o anti-Cap.~13.
      \item Vento máximo na SUPERFÍCIE — não no jato.
    \end{itemize}
    \column{0.35\textwidth}
    \figlivroslide{fig_16_1b.jpg}{0.44\textheight}{Fig.~16.1 --- o
      disco e o olho vistos do satélite}
    \column{0.35\textwidth}
    \figlivroslide{fig_16_5.jpg}{0.44\textheight}{Fig.~16.5 --- corte
      idealizado: bandas, parede, olho}
  \end{columns}
\end{frame}

\begin{frame}{Onde nascem: as condições de gênese}
  \begin{columns}
    \column{0.4\textwidth}
    \begin{itemize}
      \item TSM $\ge 26{,}5\,°$C em $\ge$50\,m de profundidade.
      \item $|\phi| \gtrsim 5°$: precisa de $f$ (Cap.~10).
      \item Cisalhamento FRACO (o contrário do Cap.~14!) $+$ umidade
            média $+$ semente (onda de leste).
      \item Atlântico Sul: cisalhamento alto $+$ TSM marginal —
            Catarina (2004) foi a exceção instrutiva.
    \end{itemize}
    \column{0.6\textwidth}
    \figlivroslide{fig_16_7.jpg}{0.5\textheight}{Fig.~16.7 ---
      trajetórias por bacia: o mapa das condições}
  \end{columns}
\end{frame}

% ---------------------------------------------------------------
\section{A máquina de Carnot}

\begin{frame}{O ciclo desenhado no céu}
  \begin{columns}
    \column{0.4\textwidth}
    \eqdestaque{$v_{max}^2 = \dfrac{C_k}{C_d}\,
      \dfrac{T_s - T_{out}}{T_{out}}\,\Delta k$}
    \vspace{0.4em}
    \begin{itemize}
      \item Entrada isotérmica no mar, subida úmida na parede, saída
            fria no topo: Carnot (T1).
      \item O limiar dos 26{,}5\,°C é Clausius--Clapeyron (Cap.~4).
      \item WISHE: vento $\to$ evaporação $\to$ núcleo quente $\to$
            vento.
    \end{itemize}
    \column{0.6\textwidth}
    \figlivroslide{fig_16_30b.jpg}{0.54\textheight}{Fig.~16.30 --- o
      ciclo in-up-out no emagrama: um motor térmico}
  \end{columns}
\end{frame}

\begin{frame}{O núcleo quente sustenta a baixa}
  \begin{columns}
    \column{0.4\textwidth}
    \begin{itemize}
      \item Anomalia de $+$10\,K concentrada no olho/parede.
      \item Hipsométrica (Cap.~1): coluna quente $=$ leve $\Rightarrow$
            $\Delta P \sim 40$--100\,hPa (T3).
      \item Contraste com o Cap.~13: lá núcleo FRIO, baixa que
            \emph{cresce} com a altura; aqui o oposto.
    \end{itemize}
    \column{0.6\textwidth}
    \figlivroslide{fig_16_28.jpg}{0.54\textheight}{Fig.~16.28 --- a
      anomalia quente em corte: o coração do furacão}
  \end{columns}
\end{frame}

\begin{frame}{Perfis radiais: Holland e o \texttt{tcpyPI}}
  \begin{columns}
    \column{0.34\textwidth}
    \begin{itemize}
      \item $P(r) = P_c + \Delta P e^{-(R_w/r)^B}$; vento gradiente
            máximo na parede (T2; Caso~2).
      \item Modelo × observações (Anita, Kate): a física simples
            acerta o perfil.
      \item MPI de sondagem real: Caso~3 (\texttt{tcpyPI}).
    \end{itemize}
    \column{0.33\textwidth}
    \figlivroslide{fig_16_31.jpg}{0.44\textheight}{Fig.~16.31 ---
      $\Delta P(r)$: modelo vs.\ furacões reais}
    \column{0.33\textwidth}
    \figlivroslide{fig_16_33.jpg}{0.44\textheight}{Fig.~16.33 ---
      $M(r)$: calmaria, pico na parede, decaimento}
  \end{columns}
\end{frame}

% ---------------------------------------------------------------
\section{Movimento}

\begin{frame}{Trajetórias: alísios, beta e recurvamento}
  \begin{columns}
    \column{0.4\textwidth}
    \begin{itemize}
      \item Fase tropical: alísios empurram para oeste (Cap.~11).
      \item Deriva beta: 1--3\,m/s para oeste E para o polo (T5).
      \item Nas latitudes médias os oestes capturam: o ``C'' clássico.
      \item Assimetria de translação: quadrantes fortes/fracos
            (Fig.~16.35b — no HS, à esquerda do deslocamento).
    \end{itemize}
    \column{0.6\textwidth}
    \figlivroslide{fig_16_21a.jpg}{0.5\textheight}{Fig.~16.21 --- as
      altas subtropicais dirigindo as trajetórias}
  \end{columns}
\end{frame}

% ---------------------------------------------------------------
\section{Perigos: a água mata mais}

\begin{frame}{Maré de tempestade}
  \begin{columns}
    \column{0.34\textwidth}
    \eqdestaque{$\dfrac{d\eta}{dx} = \dfrac{\tau}{\rho_{agua}gh}$,
      \quad $\tau = \rho_{ar}C_DV^2$}
    \vspace{0.4em}
    \begin{itemize}
      \item Vento empilha o mar $+$ barômetro invertido
            ($\sim$1\,cm/hPa, T4).
      \item Plataforma rasa e larga amplifica (Caso~4; N4).
    \end{itemize}
    \column{0.33\textwidth}
    \figlivroslide{fig_16_43a.jpg}{0.42\textheight}{Fig.~16.43 ---
      transporte de Ekman empilhando água na costa}
    \column{0.33\textwidth}
    \figlivroslide{fig_16_47.jpg}{0.42\textheight}{Fig.~16.47 ---
      surge $+$ maré $+$ ondas: o nível total}
  \end{columns}
\end{frame}

\begin{frame}{O caso Katrina em um mapa}
  \begin{columns}
    \column{0.4\textwidth}
    \begin{itemize}
      \item Surge de 3--8\,m no quadrante forte, canalizado pelo
            delta raso.
      \item Ranking de letalidade: 1º água (surge $+$ chuva), 2º
            vento, 3º tornados embutidos.
      \item Chuva não escala com a categoria: sistemas lentos e
            fracos podem ser os piores.
    \end{itemize}
    \column{0.6\textwidth}
    \figlivroslide{fig_16_41.jpg}{0.52\textheight}{Fig.~16.41 --- a
      maré de tempestade do Katrina (2005)}
  \end{columns}
\end{frame}

\begin{frame}{Temporada e clima futuro}
  \begin{columns}
    \column{0.4\textwidth}
    \begin{itemize}
      \item Pico no fim do verão: o oceano é lento para aquecer
            (Cap.~3!).
      \item $+2$\,K: MPI sobe poucos \%/K (N3), mas chuva
            ($+7$\%/K, Cap.~4) e nível do mar compõem o risco
            (Cap.~22).
      \item Em terra o motor desliga em $\sim$1 dia — a chuva não.
    \end{itemize}
    \column{0.6\textwidth}
    \figlivroslide{fig_16_40.jpg}{0.5\textheight}{Fig.~16.40 --- a
      temporada de furacões do Atlântico}
  \end{columns}
\end{frame}

% ---------------------------------------------------------------
\section{Síntese e laboratório}

\begin{frame}{O mapa do capítulo}
  \eqdestaque{$v_{max}^2 = \tfrac{C_k}{C_d}\tfrac{T_s-T_{out}}{T_{out}}
    \Delta k$;\quad $P(r)$ de Holland;\quad
    $d\eta/dx = \tau/\rho gh$;\quad gênese: 26,5\,°C, $|\phi|>5°$,
    shear fraco}
  \vspace{0.6em}
  Núcleo quente, motor de Carnot, dirigido pelos alísios, morto em
  terra — e a água como o perigo número um.
\end{frame}

\begin{frame}{Laboratório numérico --- notebook do Cap.~16}
  \texttt{notebooks/cap16\_furacoes.ipynb}
  \begin{enumerate}
    \item Carnot $+$ Clausius--Clapeyron: MPI(TSM).
    \item Holland: $P(r)$ e $v(r)$ por categoria.
    \item \texttt{tcpyPI}: intensidade potencial de sondagem real.
    \item Maré de tempestade: plataforma rasa × profunda.
  \end{enumerate}
  \vspace{0.4em}
  \textbf{Exercícios}: T1--T5 e N1--N4 nas notas; células reservadas no
  notebook.
  \vfill
  \atribuicao
\end{frame}

\end{document}
